\section{答案速查}
\noindent 本表按题号列出计算结果；条件、符号限制、拼接解与证明见对应正文。
\begingroup\small
\noindent\textbf{\hyperref[q:1.1.1]{习题 1.1.1}}\quad{\small 第\pageref{q:1.1.1}页。}
\[\begin{gathered}\dfrac{\dd x}{\dd y}=\dfrac{k}{v_0}y(a-y),\qquad
x-x_0=\dfrac{k}{v_0}\left(\dfrac{ay^2}{2}-\dfrac{y^3}{3}\right).\end{gathered}\]
\noindent\textbf{\hyperref[q:1.1.2]{习题 1.1.2}}\quad{\small 第\pageref{q:1.1.2}页。}
\[\begin{gathered}OP=PT:\ xy'+y=0;\qquad OP=OT:\ (x\mp\sqrt{x^2+y^2})y'=y;\\
OT=PT:\ (x^2-y^2)y'=2xy.\end{gathered}\]
\noindent\textbf{\hyperref[q:1.2.1]{习题 1.2.1}}\quad{\small 第\pageref{q:1.2.1}页。}
\[\begin{gathered}x(t)=2\e^{-3t}+2t+1.\end{gathered}\]
\noindent\textbf{\hyperref[q:1.2.2]{习题 1.2.2}}\quad{\small 第\pageref{q:1.2.2}页。}
\[\begin{gathered}R(t)=R_0\e^{-k(t-t_0)}.\end{gathered}\]
\noindent\textbf{\hyperref[q:1.2.3]{习题 1.2.3}}\quad{\small 第\pageref{q:1.2.3}页。}
\[\begin{gathered}v=V\tan(\theta-\omega t),\quad 0\le t\le t_*;\qquad t_*=\dfrac1\omega\arctan\dfrac{kv_0}{\sqrt{mg}},\\
s_*=s_0+\dfrac m{2k^2}\ln\left(1+\dfrac{k^2v_0^2}{mg}\right),\\
v=-V\tanh[\omega(t-t_*)],\quad
s=s_*-\dfrac m{k^2}\ln\cosh[\omega(t-t_*)]\quad(t\ge t_*).\end{gathered}\]
\noindent\textbf{\hyperref[q:1.2.4]{习题 1.2.4}}\quad{\small 第\pageref{q:1.2.4}页。}
\[\begin{gathered}\text{单摆：}\ \begin{cases}u'=v,\\v'=-(g/l)\sin u;\end{cases}\qquad
\text{弹簧：}\ \begin{cases}u'=v,\\v'=-(k/m)u-(\mu/m)v.\end{cases}\end{gathered}\]
\noindent\textbf{\hyperref[q:1.2.5(1)]{习题 1.2.5(1)}}\quad{\small 第\pageref{q:1.2.5(1)}页。}
\[\begin{gathered}x(t;0,0)=0,\qquad x(t;0,-1)=-\e^{-t}.\end{gathered}\]
\noindent\textbf{\hyperref[q:1.2.5(2)]{习题 1.2.5(2)}}\quad{\small 第\pageref{q:1.2.5(2)}页。}
\[\begin{gathered}x=-\dfrac{u'}u,\quad u''+t^2u=0;\qquad
(t_0,x_0)=(0,0),(0,-1/2),(\sqrt2,0).\end{gathered}\]
\noindent\textbf{\hyperref[q:1.2.5(3)]{习题 1.2.5(3)}}\quad{\small 第\pageref{q:1.2.5(3)}页。}
\[\begin{gathered}x=\dfrac{u'}u,\quad u''=t^2u,\quad u(0)=1,\quad
u'(0)=0\ \text{或}\ 1.\end{gathered}\]
\noindent\textbf{\hyperref[q:1.3.1(1)]{习题 1.3.1(1)}}\quad{\small 第\pageref{q:1.3.1(1)}页。}
\[\begin{gathered}x=tx'+(x')^2.\end{gathered}\]
\noindent\textbf{\hyperref[q:1.3.1(2)]{习题 1.3.1(2)}}\quad{\small 第\pageref{q:1.3.1(2)}页。}
\[\begin{gathered}x''-2x'+2x=0.\end{gathered}\]
\noindent\textbf{\hyperref[q:1.3.1(3)]{习题 1.3.1(3)}}\quad{\small 第\pageref{q:1.3.1(3)}页。}
\[\begin{gathered}(x'')^2=[1+(x')^2]^3.\end{gathered}\]
\noindent\textbf{\hyperref[q:1.3.2]{习题 1.3.2}}\quad{\small 第\pageref{q:1.3.2}页。}
\[\begin{gathered}\dfrac{\dd y}{\dd x}=\dfrac{B_y}{B_x}\ (B_x\ne0),\qquad
\dfrac{x+a}{r_+}-\dfrac{x-a}{r_-}=C.\end{gathered}\]
\noindent\textbf{\hyperref[q:2.1.2(1)]{习题 2.1.2(1)}}\quad{\small 第\pageref{q:2.1.2(1)}页。}
\[\begin{gathered}x^2+y^2=C,\quad y\ne0.\end{gathered}\]
\noindent\textbf{\hyperref[q:2.1.2(2)]{习题 2.1.2(2)}}\quad{\small 第\pageref{q:2.1.2(2)}页。}
\[\begin{gathered}y=-\dfrac1{\sin x+C},\qquad y=0.\end{gathered}\]
\noindent\textbf{\hyperref[q:2.1.2(3)]{习题 2.1.2(3)}}\quad{\small 第\pageref{q:2.1.2(3)}页。}
\[\begin{gathered}y=C\e^{x^2}.\end{gathered}\]
\noindent\textbf{\hyperref[q:2.1.2(4)]{习题 2.1.2(4)}}\quad{\small 第\pageref{q:2.1.2(4)}页。}
\[\begin{gathered}(1+x^2)(1+y^2)=Cx^2,\quad x y\ne0.\end{gathered}\]
\noindent\textbf{\hyperref[q:2.1.2(5)]{习题 2.1.2(5)}}\quad{\small 第\pageref{q:2.1.2(5)}页。}
\[\begin{gathered}\e^{-2y}=\dfrac12\e^{-4x}+C,\qquad
 y=-\dfrac12\ln\left(\dfrac12\e^{-4x}+C\right).\end{gathered}\]
\noindent\textbf{\hyperref[q:2.1.2(6)]{习题 2.1.2(6)}}\quad{\small 第\pageref{q:2.1.2(6)}页。}
\[\begin{gathered}y=\ln\left|\dfrac{x-1}{x+1}\right|+C,\quad x\ne\pm1.\end{gathered}\]
\noindent\textbf{\hyperref[q:2.1.2(7)]{习题 2.1.2(7)}}\quad{\small 第\pageref{q:2.1.2(7)}页。}
\[\begin{gathered}y=\dfrac{1+C\e^x}{1-C\e^x},\qquad y=-1.\end{gathered}\]
\noindent\textbf{\hyperref[q:2.1.2(8)]{习题 2.1.2(8)}}\quad{\small 第\pageref{q:2.1.2(8)}页。}
\[\begin{gathered}\tan y=\e^x+C,\qquad y=\dfrac\pi2+k\pi\ (k\in\mathbb Z).\end{gathered}\]
\noindent\textbf{\hyperref[q:2.1.2(9)]{习题 2.1.2(9)}}\quad{\small 第\pageref{q:2.1.2(9)}页。}
\[\begin{gathered}y=\sin(\arcsin x+C),\quad -\pi/2<\arcsin x+C<\pi/2;\qquad y=\pm1.\end{gathered}\]
\noindent\textbf{\hyperref[q:2.1.2(10)]{习题 2.1.2(10)}}\quad{\small 第\pageref{q:2.1.2(10)}页。}
\[\begin{gathered}y=-\dfrac1{x^2+C},\qquad y=0.\end{gathered}\]
\noindent\textbf{\hyperref[q:2.1.2(11)]{习题 2.1.2(11)}}\quad{\small 第\pageref{q:2.1.2(11)}页。}
\[\begin{gathered}y=C(x+1)^2.\end{gathered}\]
\noindent\textbf{\hyperref[q:2.1.2(12)]{习题 2.1.2(12)}}\quad{\small 第\pageref{q:2.1.2(12)}页。}
\[\begin{gathered}y^3+\e^y=\sin x+C.\end{gathered}\]
\noindent\textbf{\hyperref[q:2.1.2(13)]{习题 2.1.2(13)}}\quad{\small 第\pageref{q:2.1.2(13)}页。}
\[\begin{gathered}(x^2-1)(y^2-1)=C.\end{gathered}\]
\noindent\textbf{\hyperref[q:2.1.3(1)]{习题 2.1.3(1)}}\quad{\small 第\pageref{q:2.1.3(1)}页。}
\[\begin{gathered}\sin(y/x)=Cx,\quad x\ne0,\ \cos(y/x)\ne0.\end{gathered}\]
\noindent\textbf{\hyperref[q:2.1.3(2)]{习题 2.1.3(2)}}\quad{\small 第\pageref{q:2.1.3(2)}页。}
\[\begin{gathered}(x-1)^2-2(x-1)(y-2)-(y-2)^2=C,\quad x+y\ne3.\end{gathered}\]
\noindent\textbf{\hyperref[q:2.1.3(3)]{习题 2.1.3(3)}}\quad{\small 第\pageref{q:2.1.3(3)}页。}
\[\begin{gathered}\dfrac yx-\ln|y|=C\quad(xy\ne0),\qquad y=0.\end{gathered}\]
\noindent\textbf{\hyperref[q:2.1.3(4)]{习题 2.1.3(4)}}\quad{\small 第\pageref{q:2.1.3(4)}页。}
\[\begin{gathered}x+y+1=C\e^y,\qquad x+y\ne0.\end{gathered}\]
\noindent\textbf{\hyperref[q:2.1.3(5)]{习题 2.1.3(5)}}\quad{\small 第\pageref{q:2.1.3(5)}页。}
\[\begin{gathered}(x-3)^2+(x+y-1)^2=C.\end{gathered}\]
\noindent\textbf{\hyperref[q:2.1.3(6)]{习题 2.1.3(6)}}\quad{\small 第\pageref{q:2.1.3(6)}页。}
\[\begin{gathered}x^3-2y^3=Cx\quad(x\ne0).\end{gathered}\]
\noindent\textbf{\hyperref[q:2.1.3(7)]{习题 2.1.3(7)}}\quad{\small 第\pageref{q:2.1.3(7)}页。}
\[\begin{gathered}(x-y)^2+2x=C,\quad x\ne y.\end{gathered}\]
\noindent\textbf{\hyperref[q:2.1.3(8)]{习题 2.1.3(8)}}\quad{\small 第\pageref{q:2.1.3(8)}页。}
\[\begin{gathered}\dfrac{y^2}{x^2}=2\ln|x|+C,\quad xy\ne0.\end{gathered}\]
\noindent\textbf{\hyperref[q:2.1.3(9)]{习题 2.1.3(9)}}\quad{\small 第\pageref{q:2.1.3(9)}页。}
\[\begin{gathered}(x-y-2)^2+14x=C,\quad x-y\ne2.\end{gathered}\]
\noindent\textbf{\hyperref[q:2.1.3(10)]{习题 2.1.3(10)}}\quad{\small 第\pageref{q:2.1.3(10)}页。}
\[\begin{gathered}\ln|y+2|+2\arctan\dfrac{y+2}{x-3}=C,\qquad y=-2.\end{gathered}\]
\noindent\textbf{\hyperref[q:2.1.3(11)]{习题 2.1.3(11)}}\quad{\small 第\pageref{q:2.1.3(11)}页。}
\[\begin{gathered}y=x(\ln|x|+C)^2,\quad \ln|x|+C\ge0,\ x\ne0;\qquad y=0.\end{gathered}\]
\noindent\textbf{\hyperref[q:2.1.3(12)]{习题 2.1.3(12)}}\quad{\small 第\pageref{q:2.1.3(12)}页。}
\[\begin{gathered}\dfrac12\ln[(x+3)^2+(y+1)^2]+\arctan\dfrac{y+1}{x+3}=C.\end{gathered}\]
\noindent\textbf{\hyperref[q:2.1.3(13)]{习题 2.1.3(13)}}\quad{\small 第\pageref{q:2.1.3(13)}页。}
\[\begin{gathered}y=x+\dfrac{2x}{C-\ln|x|}\quad(x\ne0),\qquad y=x.\end{gathered}\]
\noindent\textbf{\hyperref[q:2.1.3(14)]{习题 2.1.3(14)}}\quad{\small 第\pageref{q:2.1.3(14)}页。}
\[\begin{gathered}-\dfrac{\cos(x+y+1)}{1+\sin(x+y+1)}=x+C;\qquad
 y=-x-1-\dfrac\pi2+2k\pi.\end{gathered}\]
\noindent\textbf{\hyperref[q:2.1.3(15)]{习题 2.1.3(15)}}\quad{\small 第\pageref{q:2.1.3(15)}页。}
\[\begin{gathered}\arcsin(y/x)=\operatorname{sgn}(x)\ln|x|+C,\qquad y=\pm x.\end{gathered}\]
\noindent\textbf{\hyperref[q:2.1.4(1)]{习题 2.1.4(1)}}\quad{\small 第\pageref{q:2.1.4(1)}页。}
\[\begin{gathered}y=\dfrac23(x+1)^{7/2}+C(x+1)^2.\end{gathered}\]
\noindent\textbf{\hyperref[q:2.1.4(2)]{习题 2.1.4(2)}}\quad{\small 第\pageref{q:2.1.4(2)}页。}
\[\begin{gathered}y=\dfrac{C-\cos x}{x},\quad x\ne0.\end{gathered}\]
\noindent\textbf{\hyperref[q:2.1.4(3)]{习题 2.1.4(3)}}\quad{\small 第\pageref{q:2.1.4(3)}页。}
\[\begin{gathered}y=\dfrac{\e^x+C}{x},\quad x>0.\end{gathered}\]
\noindent\textbf{\hyperref[q:2.1.4(4)]{习题 2.1.4(4)}}\quad{\small 第\pageref{q:2.1.4(4)}页。}
\[\begin{gathered}y=-\dfrac13+\dfrac2x+\dfrac C{x^3}.\end{gathered}\]
\noindent\textbf{\hyperref[q:2.1.4(5)]{习题 2.1.4(5)}}\quad{\small 第\pageref{q:2.1.4(5)}页。}
\[\begin{gathered}y=(x+1)[\ln(x+1)+C]+1.\end{gathered}\]
\noindent\textbf{\hyperref[q:2.1.4(6)]{习题 2.1.4(6)}}\quad{\small 第\pageref{q:2.1.4(6)}页。}
\[\begin{gathered}y=C\e^{x^2}-\dfrac12.\end{gathered}\]
\noindent\textbf{\hyperref[q:2.1.5(1)]{习题 2.1.5(1)}}\quad{\small 第\pageref{q:2.1.5(1)}页。}
\[\begin{gathered}x=y^2(C-\ln|y|),\qquad y=0.\end{gathered}\]
\noindent\textbf{\hyperref[q:2.1.5(2)]{习题 2.1.5(2)}}\quad{\small 第\pageref{q:2.1.5(2)}页。}
\[\begin{gathered}xy-\dfrac{y^4}{4}=C.\end{gathered}\]
\noindent\textbf{\hyperref[q:2.1.5(3)]{习题 2.1.5(3)}}\quad{\small 第\pageref{q:2.1.5(3)}页。}
\[\begin{gathered}x=\dfrac12\left(Ay^2-\dfrac1A\right),\quad A>0;\qquad y=0.\end{gathered}\]
\noindent\textbf{\hyperref[q:2.1.6(1)]{习题 2.1.6(1)}}\quad{\small 第\pageref{q:2.1.6(1)}页。}
\[\begin{gathered}y+\e^y=\sin x+2+\e^3.\end{gathered}\]
\noindent\textbf{\hyperref[q:2.1.6(2)]{习题 2.1.6(2)}}\quad{\small 第\pageref{q:2.1.6(2)}页。}
\[\begin{gathered}y=-\sqrt{1+2x^2},\quad x\in\R.\end{gathered}\]
\noindent\textbf{\hyperref[q:2.1.6(3)]{习题 2.1.6(3)}}\quad{\small 第\pageref{q:2.1.6(3)}页。}
\[\begin{gathered}y=\sqrt{1+2\ln\dfrac{1+\e^x}{2}},\quad
 x>\ln(2\e^{-1/2}-1).\end{gathered}\]
\noindent\textbf{\hyperref[q:2.1.6(4)]{习题 2.1.6(4)}}\quad{\small 第\pageref{q:2.1.6(4)}页。}
\[\begin{gathered}y=\dfrac32\e^{\cos x-1}.\end{gathered}\]
\noindent\textbf{\hyperref[q:2.1.6(5)]{习题 2.1.6(5)}}\quad{\small 第\pageref{q:2.1.6(5)}页。}
\[\begin{gathered}y=\ln\left(\dfrac{x^2}{2}+\dfrac{x^4}{4}+\e-\dfrac34\right),\quad x\in\R.\end{gathered}\]
\noindent\textbf{\hyperref[q:2.1.6(6)]{习题 2.1.6(6)}}\quad{\small 第\pageref{q:2.1.6(6)}页。}
\[\begin{gathered}y=\dfrac12+\dfrac32\e^{1-x^2}.\end{gathered}\]
\noindent\textbf{\hyperref[q:2.1.7(1)]{习题 2.1.7(1)}}\quad{\small 第\pageref{q:2.1.7(1)}页。}
\[\begin{gathered}y=\dfrac{8x^6}{x^8+C}\ (x\ne0),\qquad y=0.\end{gathered}\]
\noindent\textbf{\hyperref[q:2.1.7(2)]{习题 2.1.7(2)}}\quad{\small 第\pageref{q:2.1.7(2)}页。}
\[\begin{gathered}y=\dfrac1{x[C-\dfrac12(\ln x)^2]}\ (x>0),\qquad y=0.\end{gathered}\]
\noindent\textbf{\hyperref[q:2.1.7(3)]{习题 2.1.7(3)}}\quad{\small 第\pageref{q:2.1.7(3)}页。}
\[\begin{gathered}y=x^4(\ln|x|+C)^2,\quad \ln|x|+C>0.\end{gathered}\]
\noindent\textbf{\hyperref[q:2.1.7(4)]{习题 2.1.7(4)}}\quad{\small 第\pageref{q:2.1.7(4)}页。}
\[\begin{gathered}y=\dfrac1{1-x^2+C\e^{-x^2}},\qquad y=0.\end{gathered}\]
\noindent\textbf{\hyperref[q:2.1.8]{习题 2.1.8}}\quad{\small 第\pageref{q:2.1.8}页。}
\[\begin{gathered}y=\dfrac{x}{\sqrt{1+x^2}},\quad x\in\R.\end{gathered}\]
\noindent\textbf{\hyperref[q:2.1.9]{习题 2.1.9}}\quad{\small 第\pageref{q:2.1.9}页。}
\[\begin{gathered}y=\dfrac{24}{x},\quad x>0.\end{gathered}\]
\noindent\textbf{\hyperref[q:2.2.1(1)]{习题 2.2.1(1)}}\quad{\small 第\pageref{q:2.2.1(1)}页。}
\[\begin{gathered}x^3+3x^2y^2+y^4=C.\end{gathered}\]
\noindent\textbf{\hyperref[q:2.2.1(2)]{习题 2.2.1(2)}}\quad{\small 第\pageref{q:2.2.1(2)}页。}
\[\begin{gathered}\sin x+\dfrac xy+\ln|y|=C,\quad y\ne0.\end{gathered}\]
\noindent\textbf{\hyperref[q:2.2.1(3)]{习题 2.2.1(3)}}\quad{\small 第\pageref{q:2.2.1(3)}页。}
\[\begin{gathered}x^5+\dfrac32x^2y^2-xy^3+\dfrac{y^3}{3}=C.\end{gathered}\]
\noindent\textbf{\hyperref[q:2.2.1(4)]{习题 2.2.1(4)}}\quad{\small 第\pageref{q:2.2.1(4)}页。}
\[\begin{gathered}\dfrac{x^3}{3}+xy+\dfrac{y^2}{2}=C.\end{gathered}\]
\noindent\textbf{\hyperref[q:2.2.1(5)]{习题 2.2.1(5)}}\quad{\small 第\pageref{q:2.2.1(5)}页。}
\[\begin{gathered}3x^2+xy+2x+4y^2-3y=C,\quad x+8y-3\ne0.\end{gathered}\]
\noindent\textbf{\hyperref[q:2.2.1(6)]{习题 2.2.1(6)}}\quad{\small 第\pageref{q:2.2.1(6)}页。}
\[\begin{gathered}\dfrac{x^2y}{2}+\ln|y|=C,\quad y\ne0.\end{gathered}\]
\noindent\textbf{\hyperref[q:2.2.1(7)]{习题 2.2.1(7)}}\quad{\small 第\pageref{q:2.2.1(7)}页。}
\[\begin{gathered}x^3-2xy^2+y=C.\end{gathered}\]
\noindent\textbf{\hyperref[q:2.2.1(8)]{习题 2.2.1(8)}}\quad{\small 第\pageref{q:2.2.1(8)}页。}
\[\begin{gathered}\dfrac14(x^2+y^2)^2=C.\end{gathered}\]
\noindent\textbf{\hyperref[q:2.2.1(9)]{习题 2.2.1(9)}}\quad{\small 第\pageref{q:2.2.1(9)}页。}
\[\begin{gathered}3xy+\e^x+\sin y=C.\end{gathered}\]
\noindent\textbf{\hyperref[q:2.2.2(1)]{习题 2.2.2(1)}}\quad{\small 第\pageref{q:2.2.2(1)}页。}
\[\begin{gathered}x^3y+4x^2y^2+4y^3=28.\end{gathered}\]
\noindent\textbf{\hyperref[q:2.2.2(2)]{习题 2.2.2(2)}}\quad{\small 第\pageref{q:2.2.2(2)}页。}
\[\begin{gathered}x^4\e^{x+y}+x^2+y^2=1.\end{gathered}\]
\noindent\textbf{\hyperref[q:2.2.4(1)]{习题 2.2.4(1)}}\quad{\small 第\pageref{q:2.2.4(1)}页。}
\[\begin{gathered}\dfrac xy+\ln|y|=C,\qquad y=0.\end{gathered}\]
\noindent\textbf{\hyperref[q:2.2.4(2)]{习题 2.2.4(2)}}\quad{\small 第\pageref{q:2.2.4(2)}页。}
\[\begin{gathered}\sqrt{x^2+y^2}-x=C>0,\quad y>0.\end{gathered}\]
\noindent\textbf{\hyperref[q:2.2.4(3)]{习题 2.2.4(3)}}\quad{\small 第\pageref{q:2.2.4(3)}页。}
\[\begin{gathered}\arctan(y/x)-x=C\quad\text{（局部连续角度）}.\end{gathered}\]
\noindent\textbf{\hyperref[q:2.2.4(4)]{习题 2.2.4(4)}}\quad{\small 第\pageref{q:2.2.4(4)}页。}
\[\begin{gathered}\e^x\left(x^2y+\dfrac{y^3}{3}\right)=C.\end{gathered}\]
\noindent\textbf{\hyperref[q:2.2.4(5)]{习题 2.2.4(5)}}\quad{\small 第\pageref{q:2.2.4(5)}页。}
\[\begin{gathered}x^2\ln y+\dfrac13(1+y^2)^{3/2}=C,\quad y>0.\end{gathered}\]
\noindent\textbf{\hyperref[q:2.2.5]{习题 2.2.5}}\quad{\small 第\pageref{q:2.2.5}页。}
\[\begin{gathered}\mu(x,y)=y^{-\alpha}\e^{-(1-\alpha)\int a(x)\,\dd x}.\end{gathered}\]
\noindent\textbf{\hyperref[q:2.2.6]{习题 2.2.6}}\quad{\small 第\pageref{q:2.2.6}页。}
\[\begin{gathered}\mu=\dfrac1{P(x)N(y)},\qquad
 \int\dfrac{M(x)}{P(x)}\dd x+\int\dfrac{Q(y)}{N(y)}\dd y=C.\end{gathered}\]
\noindent\textbf{\hyperref[q:2.2.7]{习题 2.2.7}}\quad{\small 第\pageref{q:2.2.7}页。}
\[\begin{gathered}f(x)=c,\qquad -x^2\cos x+2x\sin x+2\cos x+\dfrac c2y^2=C.\end{gathered}\]
\noindent\textbf{\hyperref[q:2.2.8]{习题 2.2.8}}\quad{\small 第\pageref{q:2.2.8}页。}
\[\begin{gathered}f(x)=\dfrac x2+\dfrac cx,\qquad
 \dfrac{x^4}{4}+\dfrac{x^2y}{2}+cy=C.\end{gathered}\]
\noindent\textbf{\hyperref[q:2.2.9]{习题 2.2.9}}\quad{\small 第\pageref{q:2.2.9}页。}
\[\begin{gathered}a=1,\qquad \e^{-x}\sin y+x=C,\quad\cos y\ne0.\end{gathered}\]
\noindent\textbf{\hyperref[q:2.4.1(1)]{习题 2.4.1(1)}}\quad{\small 第\pageref{q:2.4.1(1)}页。}
\[\begin{gathered}x=p^2+2p,\qquad y=\dfrac23p^3+p^2+C.\end{gathered}\]
\noindent\textbf{\hyperref[q:2.4.1(2)]{习题 2.4.1(2)}}\quad{\small 第\pageref{q:2.4.1(2)}页。}
\[\begin{gathered}y=\dfrac{x^2}{2}+Cx+C^2;\qquad y=\dfrac{x^2}{4}\ \text{（奇解）}.\end{gathered}\]
\noindent\textbf{\hyperref[q:2.4.1(3)]{习题 2.4.1(3)}}\quad{\small 第\pageref{q:2.4.1(3)}页。}
\[\begin{gathered}y=x-C+\dfrac1{x-C},\quad x\ne C;\qquad y=\pm2.\end{gathered}\]
\noindent\textbf{\hyperref[q:2.4.1(4)]{习题 2.4.1(4)}}\quad{\small 第\pageref{q:2.4.1(4)}页。}
\[\begin{gathered}y=C(x-C)^2\quad(C\in\R),\qquad y=\dfrac4{27}x^3\ \text{（奇解）}.\end{gathered}\]
\noindent\textbf{\hyperref[q:2.4.1(5)]{习题 2.4.1(5)}}\quad{\small 第\pageref{q:2.4.1(5)}页。}
\[\begin{gathered}y=\dfrac{3u}{1+u^3},\quad u\ne0,-1,\\
x=C-\dfrac1u+\ln|u+1|-\dfrac12\ln(u^2-u+1)
-\sqrt3\arctan\dfrac{2u-1}{\sqrt3};\qquad y=0.\end{gathered}\]
\noindent\textbf{\hyperref[q:2.4.1(6)]{习题 2.4.1(6)}}\quad{\small 第\pageref{q:2.4.1(6)}页。}
\[\begin{gathered}y=rx+C,\qquad r\in\R,\quad r^7-5r^2+1=0.\end{gathered}\]
\noindent\textbf{\hyperref[q:2.4.2(1)]{习题 2.4.2(1)}}\quad{\small 第\pageref{q:2.4.2(1)}页。}
\[\begin{gathered}y=\sin(x+C),\qquad y=\pm1\ \text{（奇解）}.\end{gathered}\]
\noindent\textbf{\hyperref[q:2.4.2(2)]{习题 2.4.2(2)}}\quad{\small 第\pageref{q:2.4.2(2)}页。}
\[\begin{gathered}y=Cx+\dfrac1C\ (C\ne0);\qquad
 y=\pm2\sqrt x\ (x>0)\ \text{（奇解）}.\end{gathered}\]
\noindent\textbf{\hyperref[q:2.4.2(3)]{习题 2.4.2(3)}}\quad{\small 第\pageref{q:2.4.2(3)}页。}
\[\begin{gathered}y=\dfrac{(x-C)^2}{4},\qquad y=0\ \text{（奇解）}.\end{gathered}\]
\noindent\textbf{\hyperref[q:2.4.2(4)]{习题 2.4.2(4)}}\quad{\small 第\pageref{q:2.4.2(4)}页。}
\[\begin{gathered}x=C+\dfrac49p^2,\qquad y=C+\dfrac8{27}p^3;\qquad
 y=x-\dfrac4{27}\ \text{（奇解）}.\end{gathered}\]
\noindent\textbf{\hyperref[q:2.4.2(5)]{习题 2.4.2(5)}}\quad{\small 第\pageref{q:2.4.2(5)}页。}
\[\begin{gathered}y=x+\dfrac{(C-x)^2}{4}\ (x<C),\qquad y=x\ \text{（奇解）}.\end{gathered}\]
\noindent\textbf{\hyperref[q:2.4.2(6)]{习题 2.4.2(6)}}\quad{\small 第\pageref{q:2.4.2(6)}页。}
\[\begin{gathered}y=Cx+\dfrac{C^2}{2},\quad x+C>0;\qquad
 y=-\dfrac{x^2}{2}\ \text{（奇解）}.\end{gathered}\]
\noindent\textbf{\hyperref[q:2.4.3(1)]{习题 2.4.3(1)}}\quad{\small 第\pageref{q:2.4.3(1)}页。}
\[\begin{gathered}y=C_1x^5+C_2x^3+C_3x^2+C_4x+C_5,\quad x\ne0.\end{gathered}\]
\noindent\textbf{\hyperref[q:2.4.3(2)]{习题 2.4.3(2)}}\quad{\small 第\pageref{q:2.4.3(2)}页。}
\[\begin{gathered}y^2=Ax+B.\end{gathered}\]
\noindent\textbf{\hyperref[q:2.4.3(3)]{习题 2.4.3(3)}}\quad{\small 第\pageref{q:2.4.3(3)}页。}
\[\begin{gathered}y=A\left(x+\dfrac{x^3}{3}\right)+B.\end{gathered}\]
\noindent\textbf{\hyperref[q:2.4.3(4)]{习题 2.4.3(4)}}\quad{\small 第\pageref{q:2.4.3(4)}页。}
\[\begin{gathered}y=A\e^{kx},\quad A,k\in\R.\end{gathered}\]
\noindent\textbf{\hyperref[q:2.4.3(5)]{习题 2.4.3(5)}}\quad{\small 第\pageref{q:2.4.3(5)}页。}
\[\begin{gathered}y=\dfrac{(x-a)^3}{12}+b,\qquad y=b.\end{gathered}\]
\noindent\textbf{\hyperref[q:2.4.3(6)]{习题 2.4.3(6)}}\quad{\small 第\pageref{q:2.4.3(6)}页。}
\[\begin{gathered}y=1+A\e^{Bx},\quad A\ne0.\end{gathered}\]
\noindent\textbf{\hyperref[q:2.4.3(7)]{习题 2.4.3(7)}}\quad{\small 第\pageref{q:2.4.3(7)}页。}
\[\begin{gathered}y=C_1\e^{x/2}+C_2\e^{-x/2}+C_3x+C_4.\end{gathered}\]
\noindent\textbf{\hyperref[q:2.4.3(8)]{习题 2.4.3(8)}}\quad{\small 第\pageref{q:2.4.3(8)}页。}
\[\begin{gathered}y=-\dfrac xa+\dfrac{1+a^2}{a^2}\ln|1+ax|+b\quad(a\ne0),\\
y=-\dfrac{x^2}{2}+b\quad(a=0),\qquad y=\ln|x|+b.\end{gathered}\]
\noindent\textbf{\hyperref[q:2.4.3(9)]{习题 2.4.3(9)}}\quad{\small 第\pageref{q:2.4.3(9)}页。}
\[\begin{gathered}y=A\e^{kx}+B\quad(k\ne0),\qquad y=Ax+B.\end{gathered}\]
\noindent\textbf{\hyperref[q:2.4.4]{习题 2.4.4}}\quad{\small 第\pageref{q:2.4.4}页。}
\[\begin{gathered}y=1-x.\end{gathered}\]
\noindent\textbf{\hyperref[q:3.1.2]{习题 3.1.2}}\quad{\small 第\pageref{q:3.1.2}页。}
\[\begin{gathered}\det A'=48t^3-2t,\qquad (\det A)'=40t^4-3t^2+8t.\end{gathered}\]
\noindent\textbf{\hyperref[q:3.2.1]{习题 3.2.1}}\quad{\small 第\pageref{q:3.2.1}页。}
\[\begin{gathered}x=X(t)c,\qquad X(0)=I.\end{gathered}\]
\noindent\textbf{\hyperref[q:3.2.3]{习题 3.2.3}}\quad{\small 第\pageref{q:3.2.3}页。}
\[\begin{gathered}A=\begin{pmatrix}0&1\\-1&0\end{pmatrix}.\end{gathered}\]
\noindent\textbf{\hyperref[q:3.3.1]{习题 3.3.1}}\quad{\small 第\pageref{q:3.3.1}页。}
\[\begin{gathered}x(t)=\binom{-\cos t-2\sin t}{-\sin t+2\cos t}.\end{gathered}\]
\noindent\textbf{\hyperref[q:3.4.3]{习题 3.4.3}}\quad{\small 第\pageref{q:3.4.3}页。}
\[\begin{gathered}x=\dfrac{C_1\cos t+C_2\sin t}{t},\quad t\ne0.\end{gathered}\]
\noindent\textbf{\hyperref[q:3.4.4]{习题 3.4.4}}\quad{\small 第\pageref{q:3.4.4}页。}
\[\begin{gathered}x=x_1(t)\left[C_1+C_2\int\dfrac{\e^{-\int a_1(t)\dd t}}{x_1(t)^2}\dd t\right].\end{gathered}\]
\noindent\textbf{\hyperref[q:3.4.5]{习题 3.4.5}}\quad{\small 第\pageref{q:3.4.5}页。}
\[\begin{gathered}x(t)=C_1x_1(t)+C_2x_2(t)+
\int_{t_0}^t\dfrac{x_1(s)x_2(t)-x_1(t)x_2(s)}{W(s)}f(s)\dd s.\end{gathered}\]
\noindent\textbf{\hyperref[q:3.4.6]{习题 3.4.6}}\quad{\small 第\pageref{q:3.4.6}页。}
\[\begin{gathered}x=C_1\cos2t+C_2\sin2t-\dfrac{t^2}{8}\cos2t+\dfrac t{16}\sin2t.\end{gathered}\]
\noindent\textbf{\hyperref[q:3.5.1]{习题 3.5.1}}\quad{\small 第\pageref{q:3.5.1}页。}
\[\begin{gathered}x=a_0\left(1+\dfrac{t^3}{6}+\dfrac{t^6}{180}+\dfrac{t^9}{12960}+\cdots\right)
+a_1\left(t+\dfrac{t^4}{12}+\dfrac{t^7}{504}+\dfrac{t^{10}}{45360}+\cdots\right),\\
a_{3k}=\dfrac{a_0}{\prod_{j=1}^k(3j)(3j-1)},\qquad
 a_{3k+1}=\dfrac{a_1}{\prod_{j=1}^k(3j+1)(3j)}.\end{gathered}\]
\noindent\textbf{\hyperref[q:3.5.2]{习题 3.5.2}}\quad{\small 第\pageref{q:3.5.2}页。}
\[\begin{gathered}x=10+3s+15s^2+4s^3+\dfrac{75}{4}s^4+\dfrac{24}{5}s^5+\cdots,\quad s=t-1,\ |s|<1,\\
a_{2k}=10\prod_{j=1}^k\dfrac{2j+1}{2j},\qquad
 a_{2k+1}=3\prod_{j=1}^k\dfrac{2j+2}{2j+1}.\end{gathered}\]
\noindent\textbf{\hyperref[q:3.5.3]{习题 3.5.3}}\quad{\small 第\pageref{q:3.5.3}页。}
\[\begin{gathered}x=a_0\sum_{k\ge0}\left[\prod_{j=0}^{k-1}\dfrac{4j-\lambda}{(2j+2)(2j+1)}\right]t^{2k}
+a_1\sum_{k\ge0}\left[\prod_{j=0}^{k-1}\dfrac{4j+2-\lambda}{(2j+3)(2j+2)}\right]t^{2k+1},\\
\lambda=2n:\quad H_n(t)=n!\sum_{j=0}^{\lfloor n/2\rfloor}\dfrac{(-1)^j(2t)^{n-2j}}{j!(n-2j)!}.\end{gathered}\]
\noindent\textbf{\hyperref[q:4.1.1(1)]{习题 4.1.1(1)}}\quad{\small 第\pageref{q:4.1.1(1)}页。}
\[\begin{gathered}x=C_1\cos2t+C_2\sin2t.\end{gathered}\]
\noindent\textbf{\hyperref[q:4.1.1(2)]{习题 4.1.1(2)}}\quad{\small 第\pageref{q:4.1.1(2)}页。}
\[\begin{gathered}x=C_1\e^t+C_2\cos(\sqrt2t)+C_3\sin(\sqrt2t).\end{gathered}\]
\noindent\textbf{\hyperref[q:4.1.1(3)]{习题 4.1.1(3)}}\quad{\small 第\pageref{q:4.1.1(3)}页。}
\[\begin{gathered}x=\e^t(C_1\cos t+C_2\sin t)+\e^{-t}(C_3\cos t+C_4\sin t).\end{gathered}\]
\noindent\textbf{\hyperref[q:4.1.1(4)]{习题 4.1.1(4)}}\quad{\small 第\pageref{q:4.1.1(4)}页。}
\[\begin{gathered}x=(C_1+C_2t+C_3t^2)\e^t+C_4\e^{-t}.\end{gathered}\]
\noindent\textbf{\hyperref[q:4.1.1(5)]{习题 4.1.1(5)}}\quad{\small 第\pageref{q:4.1.1(5)}页。}
\[\begin{gathered}x=(C_1+C_2t)\e^{-2t}.\end{gathered}\]
\noindent\textbf{\hyperref[q:4.1.1(6)]{习题 4.1.1(6)}}\quad{\small 第\pageref{q:4.1.1(6)}页。}
\[\begin{gathered}x=\e^{-t}[C_1\cos(\sqrt3t)+C_2\sin(\sqrt3t)].\end{gathered}\]
\noindent\textbf{\hyperref[q:4.1.1(7)]{习题 4.1.1(7)}}\quad{\small 第\pageref{q:4.1.1(7)}页。}
\[\begin{gathered}x=C_1+C_2\cos2t+C_3\sin2t.\end{gathered}\]
\noindent\textbf{\hyperref[q:4.1.1(8)]{习题 4.1.1(8)}}\quad{\small 第\pageref{q:4.1.1(8)}页。}
\[\begin{gathered}x=(C_1+C_2t)\cos2t+(C_3+C_4t)\sin2t.\end{gathered}\]
\noindent\textbf{\hyperref[q:4.1.3]{习题 4.1.3}}\quad{\small 第\pageref{q:4.1.3}页。}
\[\begin{gathered}\delta<\omega:\ x=\e^{-\delta t}(C_1\cos\nu t+C_2\sin\nu t),\quad\nu=\sqrt{\omega^2-\delta^2};\\
\delta=\omega:\ x=(C_1+C_2t)\e^{-\omega t};\\
\delta>\omega:\ x=C_1\e^{(-\delta+\sqrt{\delta^2-\omega^2})t}
+C_2\e^{(-\delta-\sqrt{\delta^2-\omega^2})t}.\end{gathered}\]
\noindent\textbf{\hyperref[q:4.2.1(1)]{习题 4.2.1(1)}}\quad{\small 第\pageref{q:4.2.1(1)}页。}
\[\begin{gathered}x=C_1+C_2\e^{-t}+\dfrac{t^3}{3}-t^2+3t.\end{gathered}\]
\noindent\textbf{\hyperref[q:4.2.1(2)]{习题 4.2.1(2)}}\quad{\small 第\pageref{q:4.2.1(2)}页。}
\[\begin{gathered}x=C_1\cos2t+C_2\sin2t+\dfrac t4\sin2t-\dfrac1{12}\cos4t.\end{gathered}\]
\noindent\textbf{\hyperref[q:4.2.1(3)]{习题 4.2.1(3)}}\quad{\small 第\pageref{q:4.2.1(3)}页。}
\[\begin{gathered}x=C_1\cos2t+C_2\sin2t-\dfrac{t^2}{8}\cos2t+\dfrac t{16}\sin2t.\end{gathered}\]
\noindent\textbf{\hyperref[q:4.2.1(4)]{习题 4.2.1(4)}}\quad{\small 第\pageref{q:4.2.1(4)}页。}
\[\begin{gathered}x=C_1+C_2\e^t+C_3\e^{3t}+\dfrac{t^3}{9}+\dfrac{4t^2}{9}+\dfrac{26t}{27}.\end{gathered}\]
\noindent\textbf{\hyperref[q:4.2.1(5)]{习题 4.2.1(5)}}\quad{\small 第\pageref{q:4.2.1(5)}页。}
\[\begin{gathered}x=C_1\e^t+C_2\e^{2t}-\e^{2t}\sin(\e^{-t}).\end{gathered}\]
\noindent\textbf{\hyperref[q:4.2.1(6)]{习题 4.2.1(6)}}\quad{\small 第\pageref{q:4.2.1(6)}页。}
\[\begin{gathered}x=C_1\e^t+(C_2+C_3t)\e^{2t}+\dfrac12\e^{3t}.\end{gathered}\]
\noindent\textbf{\hyperref[q:4.2.1(7)]{习题 4.2.1(7)}}\quad{\small 第\pageref{q:4.2.1(7)}页。}
\[\begin{gathered}x=C_1\e^{\sqrt2t}+C_2\e^{-\sqrt2t}+\e^{t^2}.\end{gathered}\]
\noindent\textbf{\hyperref[q:4.2.1(8)]{习题 4.2.1(8)}}\quad{\small 第\pageref{q:4.2.1(8)}页。}
\[\begin{gathered}x=\e^t\left(C_1+C_2t+\dfrac{t^3}{6}\right).\end{gathered}\]
\noindent\textbf{\hyperref[q:4.2.2]{习题 4.2.2}}\quad{\small 第\pageref{q:4.2.2}页。}
\[\begin{gathered}a_2u_{ss}+(a_1-a_2)u_s+a_0u=f(\varepsilon\e^s),\quad
\varepsilon=\operatorname{sgn}(t).\end{gathered}\]
\noindent\textbf{\hyperref[q:4.3.1]{习题 4.3.1}}\quad{\small 第\pageref{q:4.3.1}页。}
\[\begin{gathered}x_1=C_1\e^t+C_2\cos t+C_3\sin t-t-2+\dfrac25\e^{2t},\\
x_2=2C_1\e^t+(C_2-C_3)\cos t+(C_2+C_3)\sin t-3t-3+\dfrac35\e^{2t}.\end{gathered}\]
\noindent\textbf{\hyperref[q:4.3.3(1)]{习题 4.3.3(1)}}\quad{\small 第\pageref{q:4.3.3(1)}页。}
\[\begin{gathered}\begin{pmatrix}x\\y\\z\end{pmatrix}
=C_1\e^t\begin{pmatrix}3\\2\\3\end{pmatrix}
+C_2\e^{2t}\begin{pmatrix}4\\1\\1\end{pmatrix}
+C_3\e^{3t}\begin{pmatrix}2\\2\\3\end{pmatrix}.\end{gathered}\]
\noindent\textbf{\hyperref[q:4.3.3(2)]{习题 4.3.3(2)}}\quad{\small 第\pageref{q:4.3.3(2)}页。}
\[\begin{gathered}x=C_1\e^t+C_2\e^{-t},\qquad y=C_1\e^t-C_2\e^{-t}.\end{gathered}\]
\noindent\textbf{\hyperref[q:4.3.3(3)]{习题 4.3.3(3)}}\quad{\small 第\pageref{q:4.3.3(3)}页。}
\[\begin{gathered}\begin{pmatrix}x\\y\\z\end{pmatrix}=C_1\e^{5t}v+
\e^{2t}\{C_2(a\cos t-b\sin t)+C_3(a\sin t+b\cos t)\}.\end{gathered}\]
\noindent\textbf{\hyperref[q:4.3.3(4)]{习题 4.3.3(4)}}\quad{\small 第\pageref{q:4.3.3(4)}页。}
\[\begin{gathered}\begin{pmatrix}x\\y\\z\end{pmatrix}=C_1\e^{-2t}\begin{pmatrix}0\\0\\1\end{pmatrix}
+C_2\e^{\lambda_+t}v(\lambda_+)+C_3\e^{\lambda_-t}v(\lambda_-).\end{gathered}\]
\noindent\textbf{\hyperref[q:4.3.3(5)]{习题 4.3.3(5)}}\quad{\small 第\pageref{q:4.3.3(5)}页。}
\[\begin{gathered}\begin{pmatrix}x\\y\\z\end{pmatrix}
=C_1\begin{pmatrix}25\\45\\9\end{pmatrix}
+C_2\e^{4t}\begin{pmatrix}1\\1\\1\end{pmatrix}
+C_3\e^{9t}\begin{pmatrix}1\\0\\0\end{pmatrix}.\end{gathered}\]
\noindent\textbf{\hyperref[q:4.3.3(6)]{习题 4.3.3(6)}}\quad{\small 第\pageref{q:4.3.3(6)}页。}
\[\begin{gathered}x=C_1\e^t+C_2\e^{5t},\quad y=C_1\e^t-C_2\e^{5t},\quad z=C_3\e^{5t}.\end{gathered}\]
\noindent\textbf{\hyperref[q:4.3.3(7)]{习题 4.3.3(7)}}\quad{\small 第\pageref{q:4.3.3(7)}页。}
\[\begin{gathered}x=\e^{2t}[C_1+C_2t+C_3(3t-t^2/2)],\\
y=\e^{2t}(C_2-C_3t),\qquad z=C_3\e^{2t}.\end{gathered}\]
\noindent\textbf{\hyperref[q:4.4.1]{习题 4.4.1}}\quad{\small 第\pageref{q:4.4.1}页。}
\[\begin{gathered}x(t)=\dfrac14\cos8t\ \mathrm m,\qquad T=\dfrac\pi4\ \mathrm s.\end{gathered}\]
\noindent\textbf{\hyperref[q:4.4.2]{习题 4.4.2}}\quad{\small 第\pageref{q:4.4.2}页。}
\[\begin{gathered}v=A\cos(\omega_0t)+B\sin(\omega_0t),\quad \omega_0=\dfrac1{\sqrt{LC}},\quad T=2\pi\sqrt{LC}.\end{gathered}\]
\noindent\textbf{\hyperref[q:4.4.3]{习题 4.4.3}}\quad{\small 第\pageref{q:4.4.3}页。}
\[\begin{gathered}v=A\cos\omega_0t+B\sin\omega_0t-\dfrac{E_0\omega_0}{2}t\cos\omega_0t.\end{gathered}\]
\noindent\textbf{\hyperref[q:5.1.1]{习题 5.1.1}}\quad{\small 第\pageref{q:5.1.1}页。}
\[\begin{gathered}|\varphi_n(t)-\varphi(t)|\le\dfrac{ML^nh^{n+1}}{(n+1)!},\quad |t-t_0|\le h.\end{gathered}\]
\noindent\textbf{\hyperref[q:5.1.2]{习题 5.1.2}}\quad{\small 第\pageref{q:5.1.2}页。}
\[\begin{gathered}x(t)=\e^{At}x_0.\end{gathered}\]
\noindent\textbf{\hyperref[q:5.1.3]{习题 5.1.3}}\quad{\small 第\pageref{q:5.1.3}页。}
\[\begin{gathered}\varphi_3=1+t+t^2+t^3+\dfrac{2t^4}{3}+\dfrac{t^5}{3}+\dfrac{t^6}{9}+\dfrac{t^7}{63}.\end{gathered}\]
\noindent\textbf{\hyperref[q:5.1.4]{习题 5.1.4}}\quad{\small 第\pageref{q:5.1.4}页。}
\[\begin{gathered}I=[-1/2,1/2],\qquad
\varphi_3=-\dfrac{t^2}{2}+\dfrac{t^5}{20}-\dfrac{t^8}{160}+\dfrac{t^{11}}{4400},\\
|\varphi_3-\varphi|\le\dfrac{2\cdot2^3(1/2)^4}{4!}=\dfrac1{24}.\end{gathered}\]
\noindent\textbf{\hyperref[q:5.1.5]{习题 5.1.5}}\quad{\small 第\pageref{q:5.1.5}页。}
\[\begin{gathered}I=[-1/2,1/2],\qquad
\varphi_3=\dfrac{t^2}{2}+\dfrac{t^5}{20}+\dfrac{t^8}{160}+\dfrac{t^{11}}{4400}.\end{gathered}\]
\noindent\textbf{\hyperref[q:5.1.6]{习题 5.1.6}}\quad{\small 第\pageref{q:5.1.6}页。}
\[\begin{gathered}\varphi_n=2\sum_{k=0}^n\dfrac{t^k}{k!}-t-2+\dfrac{t^{n+1}}{(n+1)!},\quad
x=\lim_{n\to\infty}\varphi_n=2\e^t-t-2.\end{gathered}\]
\noindent\textbf{\hyperref[q:5.1.7]{习题 5.1.7}}\quad{\small 第\pageref{q:5.1.7}页。}
\[\begin{gathered}x(t)=\e^{A(t)}\left[x_0+\int_{t_0}^t\e^{-A(s)}Q(s)\dd s\right].\end{gathered}\]
\noindent\textbf{\hyperref[q:5.2.1]{习题 5.2.1}}\quad{\small 第\pageref{q:5.2.1}页。}
\[\begin{gathered}X_k=(5/4)^k\ (k=0,1,\ldots,4),\qquad
 X_{-k}=(3/4)^k\ (k=0,1,\ldots,4).\end{gathered}\]
\noindent\textbf{\hyperref[q:5.3.5(1)]{习题 5.3.5(1)}}\quad{\small 第\pageref{q:5.3.5(1)}页。}
\[\begin{gathered}x=t\ln t,\quad (\alpha,\beta)=(0,\infty),\quad
\lim_{t\downarrow0}x(t)=0,\quad\lim_{t\to\infty}x(t)=+\infty.\end{gathered}\]
\noindent\textbf{\hyperref[q:5.3.5(2)]{习题 5.3.5(2)}}\quad{\small 第\pageref{q:5.3.5(2)}页。}
\[\begin{gathered}x=\dfrac{1+\e^t}{1-\e^t},\quad(\alpha,\beta)=(0,\infty),\quad
\lim_{t\downarrow0}x(t)=-\infty,\quad\lim_{t\to\infty}x(t)=-1.\end{gathered}\]
\noindent\textbf{\hyperref[q:5.4.2]{习题 5.4.2}}\quad{\small 第\pageref{q:5.4.2}页。}
\[\begin{gathered}u(t)\le\delta\e^{L(t-t_0)}+\dfrac\varepsilon L[\e^{L(t-t_0)}-1]\quad(L>0),\\
u(t)\le\delta+\varepsilon(t-t_0)\quad(L=0).\end{gathered}\]
\noindent\textbf{\hyperref[q:5.5.3]{习题 5.5.3}}\quad{\small 第\pageref{q:5.5.3}页。}
\[\begin{gathered}\left.\dfrac{\partial x}{\partial t_0}\right|_{t_0=x_0=0}=0,\qquad
\left.\dfrac{\partial x}{\partial x_0}\right|_{t_0=x_0=0}=\e^{t^2/2}.\end{gathered}\]
\noindent\textbf{\hyperref[q:5.6.2]{习题 5.6.2}}\quad{\small 第\pageref{q:5.6.2}页。}
\[\begin{gathered}X_{10}^{\mathrm E}=1.852594670,\quad X_{10}^{\mathrm H}=2.020955715,\quad X_{10}^{\mathrm {RK4}}=2.027392347.\end{gathered}\]
\noindent\textbf{\hyperref[q:5.6.3]{习题 5.6.3}}\quad{\small 第\pageref{q:5.6.3}页。}
\[\begin{gathered}I_{\mathrm{RK4}}=0.7468241838759147,\qquad I_{\mathrm{ref}}=0.746824132812427,\\ |I_{\mathrm{RK4}}-I_{\mathrm{ref}}|\approx 5.106\times10^{-8}.\end{gathered}\]
\endgroup